\subsection{Bias-Variance and Approximation-Estimation Tradeoffs}

Learning quality depends both on the richness of the model class and on the variability introduced by fitting that class from a finite sample.

\begin{proposition} \label{prop:sl-bias-variance-regression} \hlt{Bias-Variance Decomposition}\\
Let $Y=f_0(X)+\varepsilon$ with $\mathbb{E}[\varepsilon\mid X]=0$ and $\operatorname{Var}(\varepsilon\mid X)=\sigma^2(X)$.  If a learning algorithm trained on a random sample $S$ returns $\widehat{f}_S$, then at a fixed input $x$,
\[
\mathbb{E}_{S,Y\mid x}\!\left[(Y-\widehat{f}_S(x))^2\right]
=
\sigma^2(x)
+\left(f_0(x)-\mathbb{E}_S\widehat{f}_S(x)\right)^2
+\operatorname{Var}_S(\widehat{f}_S(x)).
\]
\end{proposition}

\begin{definition} \label{def:sl-approximation-estimation-tradeoff} \hlt{Approximation and Estimation Errors}\\
For a class $\mathcal{F}$ and target risk $R$, the \hlt{approximation error} is
\[
\inf_{f\in\mathcal{F}}R(f)-R^*,
\]
while the \hlt{estimation error} is the excess cost of choosing a data-dependent member of $\mathcal{F}$.  Enlarging $\mathcal{F}$ often lowers approximation error while increasing estimation difficulty.
\end{definition}

\begin{remark} \label{rem:sl-model-complexity-tradeoff} \hlt{Model Complexity}\\
Classical learning theory often predicts a U-shaped validation curve as complexity grows.  Modern overparameterized models may enter interpolation regimes where optimization, implicit regularization, and data geometry change this simple picture, so complexity must be measured by more than parameter count.
\end{remark}
